\documentclass[10pt,a4paper]{amsart}
\usepackage[margin=19mm]{geometry}
\usepackage{amsmath,amssymb,mathtools}
\usepackage[scr=rsfs,cal=euler]{mathalfa}
\usepackage{xfrac}
\usepackage[unicode,hidelinks,bookmarks=false]{hyperref}
\newcommand{\ie}{i.e.\ }
\newcommand{\cf}{cf.\ }
\newcommand{\resp}{respectively\ }
\newcommand{\Subsection}{\subsection}
\providecommand{\nobreakdash}{\nobreak\hbox{-}}
\setlength{\emergencystretch}{2em}
\multlinegap0pt
\usepackage{amssymb}
\usepackage{tikz}
\usepackage[shortlabels]{enumitem}
\newtheorem{lemma}{Lemma}
\newtheorem{proposition}{Proposition}
\newcommand\R{\mathbb{R}}
\newcommand{\abs}[1]{\left\lvert #1\right\rvert}
\newcommand{\cA}{\mathcal A}
\newcommand{\cdo}{\, \cdot \,}
\newcommand\dd{\,{\rm d}}
\newcommand{\dist}{\operatorname{dist}}
\newcommand{\ip}[2]{\left\langle #1,#2\right\rangle}
\newcommand{\norm}[2]{\left\lVert #1\right\rVert_{#2}}
\newcommand{\supp}{{\rm supp}{\hspace{.05cm}}}
\title{Stability of the Riesz-type inequalities under gluing}
\author{Dangyang He}
\date{Source: arXiv:2609.21399v1 (CC BY 4.0)}
\begin{document}
\maketitle

\noindent\emph{Source and licence.} Stability of the Riesz-type inequalities under gluing. Version arXiv:2609.21399v1 (CC BY 4.0), distributed by arXiv under the Creative Commons Attribution 4.0 International licence, http://creativecommons.org/licenses/by/4.0/, as recorded in the arXiv licence metadata for this exact version and shown on the abstract page https://arxiv.org/abs/2609.21399v1. Full licence notice: \texttt{licenses/arxiv-2609.21399-CC-BY-4.0.txt}.\par\smallskip

\noindent\textit{Editorial scope.} Counted unit: the article's proposition prop:small-pairing (source0.tex lines 1712--1721). The article's complete original proof of it is delivered with it (source0.tex lines 1723--1900, one \texttt{proof} environment of 178 lines). No mathematics is written, repaired or re-derived: the statement and the proof are the article's own lines, in the article's own notation, and no retained line is substituted or re-flowed.  Retained uncounted as the source-local context the proof needs: lines 1252--1255; lines 1258--1261; lines 1264--1269; lines 1281--1284; lines 1295--1297; lines 1301--1303; lines 1365--1369; lines 1505--1508; lines 1568--1571; lines 1634--1657.  Nothing was omitted from any retained span, so the package carries no omission record.  No retained line contains a citation, so the wrapper carries no bibliography and no bibliography entry.  No retained line contains a figure environment, an image-inclusion command or drawing code, so no image is delivered and the package declares no asset dependency.  Editorial formatting only, in addition: an AMS \texttt{amsart} preamble that restates the article's own declarations and macros used by the retained lines, namely the environments \texttt{lemma}, \texttt{proposition} on separate counters, together with the standard packages those lines need.  The retained environments print in this document as Lemma 1, Proposition 1 in order of appearance, whereas the article prints the same items with the numbering of its own full document.  The complete native source of the article as distributed in the arXiv e-print for this exact version is bundled unchanged as \texttt{sources/210956/source0.tex}.
\begin{equation}\label{eq:exterior-low-Riesz}
 \norm{\nabla L_B^{-1/2}u}{L^r(\Omega)}
 \leq C_r\norm u{L^r(\Omega)},\qquad 1<r\leq2,
\end{equation}

\begin{equation}\label{eq:whole-space-reverse}
 \norm{L^{1/2}v}{L^r(\R^n)}
 \leq C_r\norm{\nabla v}{L^r(\R^n)},\qquad 1<r<\infty,
\end{equation}

\begin{equation}\label{eq:localized-square-root}
 \norm{L_B^{1/2}v}{L^r(\Omega)}
 \leq C_r\bigl(\norm v{L^r(\Omega)}
                    +\norm{\nabla v}{L^r(\Omega)}\bigr),
 \qquad 1<r<\infty.
\end{equation}

\begin{equation}\label{eq:core-reverse}
 \norm{(1+L_{U,B})^{1/2}v}{L^r(U)}
 \leq C_r\bigl(\norm v{L^r(U)}+\norm{\nabla v}{L^r(U)}\bigr).
\end{equation}

\begin{equation}\label{eq:exterior-Sobolev}
 \norm f{L^{r^*}(\Omega)}\leq C_r\norm{\nabla f}{L^r(\Omega)}
\end{equation}

\begin{equation}\label{eq:exterior-local-Poincare}
 \norm f{L^r(K\cap\Omega)}\leq C_{K,r}\norm{\nabla f}{L^r(\Omega)}.
\end{equation}

\begin{equation}\label{eq:form-commutator}
 \ip{[L_B,\phi]v}{\psi}
 =\int_\Omega vA\nabla\phi\mathbin{\cdo}\nabla\overline\psi
 -\int_\Omega\overline\psi A\nabla\phi\mathbin{\cdo}\nabla v.
\end{equation}

\begin{equation}\label{eq:m-small}
 \abs{m_s(u)}+\abs{m_s''(u)}
 \leq Cs\norm u{L^r(\Omega)}.
\end{equation}

\begin{equation}\label{eq:domain-small-off}
 \norm{e^{-s\sqrt{L_B}}f}{W^{1,2}(K\cap\Omega)}
 \leq C_{K,r}s\norm f{L^r(\Omega)}.
\end{equation}

\begin{lemma}\label{lem:Wronskian}
For $j=1,2$, let $T_j$ be a nonnegative self-adjoint operator on a Hilbert
space $\mathcal H_j$.  Let $\theta:\mathcal H_2\to\mathcal H_1$ be bounded
and map $D(T_2^{1/2})$ continuously into $D(T_1^{1/2})$.  For
$a\in\mathcal H_1$ and $b\in\mathcal H_2$, set, for $s>0$,
\[
 x_s=e^{-sT_1^{1/2}}a,
 \qquad y_s=e^{-sT_2^{1/2}}b,
\]
and
\[
 D(s)=\ip{T_1^{1/2}x_s}{\theta y_s}
      -\ip{x_s}{\theta T_2^{1/2}y_s}.
\]
The expression on the right of the identity below is a form-dual pairing:
\[
 \ip{x}{(T_1\theta-\theta T_2)y}
 :=\mathfrak q_{T_1}(x,\theta y)-\ip{x}{\theta T_2y}.
\]
Then, in the form sense,
\begin{equation}\label{eq:Wronskian-derivative}
 D'(s)=-\ip{x_s}{(T_1\theta-\theta T_2)y_s}.
\end{equation}
\end{lemma}

\begin{proposition}\label{prop:small-pairing}
Let $1<p<2$ and let $f,u\in C_c^\infty(\overline\Omega)$.  Then
\begin{equation}\label{eq:small-pairing}
 \sup_{0<\delta<1}
 \abs{\int_\delta^1
 \ip{e^{-s\sqrt{L_N}}f}{F_{N,s}u}\,\dd s}
 \leq C\norm{\nabla f}{L^p(\Omega)}\norm u{L^{p'}(\Omega)}.
\end{equation}
Moreover, the integral has a limit as $\delta\downarrow0$.
\end{proposition}

\begin{proof}
Let $K$ contain the supports of $\nabla\phi_c$, $\nabla\phi_e$,
$\phi_c$, and $\phi_e-1$.  Choose
$\chi\in C_c^\infty(\R^n)$ equal to one on a small neighborhood of $K$,
and put
\[
 f_{\mathrm{loc}}=\chi f,\qquad f_{\mathrm{far}}=(1-\chi)f.
\]
Because $1<p<2\leq n$, we have $p<n$.  The compact Poincar\'e estimate
\eqref{eq:exterior-local-Poincare}, applied on the support of $\chi$, yields
\[
 \norm{f_{\mathrm{loc}}}{W^{1,p}(\Omega)}
 \leq C\norm{\nabla f}{L^p(\Omega)}.
\]
Indeed,
\(\nabla(\chi f)=\chi\nabla f+f\nabla\chi\), and the compact Poincar\'e
estimate controls the second term as well as \(\|\chi f\|_p\).
The global Sobolev inequality \eqref{eq:exterior-Sobolev} also implies, with
$p^*=np/(n-p)$,
\begin{equation}\label{eq:ffar-pstar}
 \norm{f_{\mathrm{far}}}{L^{p^*}(\Omega)}
 \leq C\norm{\nabla f}{L^p(\Omega)}.
\end{equation}
Here the same product calculation yields
\(\|\nabla f_{\mathrm{far}}\|_p\leq C\|\nabla f\|_p\).

We first treat the local part.  To keep the endpoint calculation readable,
write
\[
 X_s=e^{-s\sqrt{L_N}}f_{\mathrm{loc}},\qquad
 U_{c,s}=e^{-s\sqrt{H_{c,N}}}(\rho_cu),\qquad
 U_{e,s}=e^{-s\sqrt L}(\rho_eu).
\]
For the core, set
\[
 D_c(s)=\ip{L_N^{1/2}X_s}{\phi_cU_{c,s}}
       -\ip{X_s}{\phi_cH_{c,N}^{1/2}U_{c,s}}.
\]
Lemma~\ref{lem:Wronskian}, with
$T_1=L_N$, $T_2=H_{c,N}$, and $\theta=\phi_c$, yields for
$0<\delta<1$
\begin{equation}\label{eq:core-Wronskian-truncated}
 \int_\delta^1\ip{X_s}
 {[L_N,\phi_c]U_{c,s}-\phi_cU_{c,s}}\,\dd s
 =D_c(\delta)-D_c(1).
\end{equation}
At zero, \eqref{eq:localized-square-root}, \eqref{eq:core-reverse}, and
H\"older's inequality yield
\begin{align*}
 \abs{\ip{L_N^{1/2}f_{\mathrm{loc}}}{\phi_c\rho_cu}}
 &\leq C\norm{f_{\mathrm{loc}}}{W^{1,p}(\Omega)}
          \norm u{L^{p'}(\Omega)},\\
 \abs{\ip{f_{\mathrm{loc}}}
 {\phi_cH_{c,N}^{1/2}(\rho_cu)}}
 &=\abs{\ip{H_{c,N}^{1/2}(\phi_cf_{\mathrm{loc}})}{\rho_cu}}\\
 &\leq C\norm{\phi_cf_{\mathrm{loc}}}{W^{1,p}(\Omega_c)}
          \norm u{L^{p'}(\Omega)}.
\end{align*}
To pass to the lower endpoint, commute $L_N^{1/2}$ with its Poisson
semigroup in the first pairing and move $H_{c,N}^{1/2}$ onto
$\phi_cX_\delta$ in the second.  The low-exponent Riesz bound,
\eqref{eq:localized-square-root}, and strong semigroup continuity imply
$X_\delta\to f_{\mathrm{loc}}$ in $W^{1,p}$; then
\eqref{eq:core-reverse} and the strong $L^{p'}$-continuity of the core
semigroup imply $D_c(\delta)\to D_c(0)$, and
$|D_c(0)|\leq C\|\nabla f\|_p\|u\|_{p'}$.  At $s=1$, analyticity of the
two Poisson semigroups implies
\begin{align*}
 \norm{L_N^{1/2}e^{-\sqrt{L_N}}f_{\mathrm{loc}}}{L^p}
 &\leq C\norm{f_{\mathrm{loc}}}{L^p},&
 \norm{e^{-\sqrt{H_{c,N}}}(\rho_cu)}{L^{p'}}
 &\leq C\norm u{L^{p'}},\\
 \norm{e^{-\sqrt{L_N}}f_{\mathrm{loc}}}{L^p}
 &\leq\norm{f_{\mathrm{loc}}}{L^p},&
 \norm{H_{c,N}^{1/2}e^{-\sqrt{H_{c,N}}}(\rho_cu)}{L^{p'}}
 &\leq C\norm u{L^{p'}}.
\end{align*}
Consequently $|D_c(1)|\leq
C\|\nabla f\|_p\|u\|_{p'}$.

For the end, define
\[
 D_e(s)=\ip{L_N^{1/2}X_s}{\phi_eU_{e,s}}
       -\ip{X_s}{\phi_eL^{1/2}U_{e,s}}.
\]
Lemma~\ref{lem:Wronskian}, now with $T_2=L$, yields
\begin{equation}\label{eq:end-Wronskian-truncated}
 \int_\delta^1\ip{X_s}{[L_N,\phi_e]U_{e,s}}\,\dd s
 =D_e(\delta)-D_e(1).
\end{equation}
The first term of $D_e(0)$ is bounded by
\eqref{eq:localized-square-root}.  For the second term, $\phi_ef_{\mathrm{loc}}$
vanishes near $\partial\Omega$ and therefore extends by zero to a function
in $W^{1,p}(\R^n)$.  Self-adjointness and
\eqref{eq:whole-space-reverse} imply
\[
 \abs{\ip{f_{\mathrm{loc}}}{\phi_eL^{1/2}(\rho_eu)}}
 =\abs{\ip{L^{1/2}(\phi_ef_{\mathrm{loc}})}{\rho_eu}}
 \leq C\norm{\phi_ef_{\mathrm{loc}}}{W^{1,p}(\Omega)}
          \norm u{L^{p'}}.
\]
The convergence $D_e(\delta)\to D_e(0)$ follows by the same argument as
before.  The two terms of $D_e(1)$ are bounded
exactly as at the core endpoint, using
analyticity on $L^p(\Omega)$ and $L^{p'}(\R^n)$.  Hence
\[
 \abs{D_e(0)}+\abs{D_e(1)}
 \leq C\norm{\nabla f}{L^p}\norm u{L^{p'}}.
\]
Equations \eqref{eq:core-Wronskian-truncated} and
\eqref{eq:end-Wronskian-truncated} control the uncorrected local error and
show that its integral has a limit at zero.

We next estimate the two local correction terms absolutely.  Since $m_s$
is spatially constant, \eqref{eq:form-commutator} implies
\[
 \abs{\ip{[L_N,\phi_e]m_s(u)}{X_s}}
 \leq C\abs{m_s(u)}\norm{\nabla X_s}{L^p(\Omega)}.
\]
Factor $\nabla e^{-s\sqrt{L_N}}=(\nabla L_N^{-1/2})L_N^{1/2}e^{-s\sqrt{L_N}}$. The Riesz estimate \eqref{eq:exterior-low-Riesz} and
analyticity imply
\[
 \norm{\nabla X_s}{L^p}
 \leq Cs^{-1}\norm{f_{\mathrm{loc}}}{L^p},\qquad 0<s\leq1.
\]
Together with \eqref{eq:m-small}, this produces an $s$-independent integrable
bound.  Similarly, by contractivity and \eqref{eq:m-small},
\[
 \abs{\ip{(\phi_e-1)m_s''(u)}{X_s}}
 \leq Cs\norm u{L^{p'}}\norm{f_{\mathrm{loc}}}{L^p}.
\]
It remains to treat $f_{\rm far}$.  Choose bounded open sets $K_1,K_2$ with
\[
 K\subset K_1,\qquad \overline{K_1}\subset K_2,\qquad
 \dist(K_2,\supp f_{\rm far})>0.
\]
Applying the off-diagonal estimate \eqref{eq:domain-small-off} with exponent
$p^*$ and using \eqref{eq:ffar-pstar}, we obtain
\begin{equation}\label{eq:far-Poisson-small}
 \norm{e^{-s\sqrt{L_N}}f_{\rm far}}{W^{1,2}(K_1\cap\Omega)}
 \leq Cs\norm{\nabla f}{L^p},\qquad0<s\leq1.
\end{equation}
For the core model, the $L^2$ Kato estimate and spectral calculus yield
\begin{align*}
 \norm{U_{c,s}}{L^2(\Omega_c)}
 &\leq C\norm u{L^{p'}},\\
 \norm{\nabla U_{c,s}}{L^2(\Omega_c)}
 &\leq C\norm{H_{c,N}^{1/2}e^{-s\sqrt{H_{c,N}}}(\rho_cu)}{L^2}
 \leq Cs^{-1}\norm u{L^{p'}}.
\end{align*}
Here $L^{p'}(\Omega_c)\subset L^2(\Omega_c)$ because $p'>2$ and the core has finite measure.  On the end transition region, the supports of
$\rho_eu$ and $\nabla\phi_e$ are separated.  The whole-space version of
the calculation proving \eqref{eq:domain-small-off} shows
\begin{equation}\label{eq:end-small-W12}
 \norm{U_{e,s}}{W^{1,2}(\cA_1)}
 +\abs{m_s(u)}+\abs{m_s''(u)}
 \leq Cs\norm u{L^{p'}}.
\end{equation}
Let $X_s^{\mathrm{far}}=e^{-s\sqrt{L_N}}f_{\rm far}$.  The form commutator
\eqref{eq:form-commutator} and
\eqref{eq:far-Poisson-small}--\eqref{eq:end-small-W12} now conclude
\begin{align*}
 \abs{\ip{[L_N,\phi_c]U_{c,s}}{X_s^{\mathrm{far}}}}
 &\leq C\norm{U_{c,s}}{W^{1,2}}
          \norm{X_s^{\mathrm{far}}}{W^{1,2}(K_1\cap\Omega)}
 \leq C\norm u{p'}\norm{\nabla f}p,\\
 \abs{\ip{\phi_cU_{c,s}}{X_s^{\mathrm{far}}}}
 &\leq C\norm{U_{c,s}}2\norm{X_s^{\mathrm{far}}}{L^2(K_1\cap\Omega)}
 \leq Cs\norm u{p'}\norm{\nabla f}p,\\
 \abs{\ip{[L_N,\phi_e](U_{e,s}-m_s)}{X_s^{\mathrm{far}}}}
 &\leq Cs^2\norm u{p'}\norm{\nabla f}p,\\
 \abs{\ip{(\phi_e-1)m_s''}{X_s^{\mathrm{far}}}}
 &\leq Cs^2\norm u{p'}\norm{\nabla f}p.
\end{align*}
Every far-part term is absolutely $s$-integrable on $(0,1)$.  Together with
the two Wronskian identities and the local correction estimates, this proves
\eqref{eq:small-pairing} and the existence of the improper integral.
\end{proof}

\end{document}
